\section{CROMODINÁMICA CUÁNTICA (QCD)}
\vspace{-3pt}
$$\mathcal{L}_{QCD} = \sum_{q} \bar{q}(i\cancel{D}-m_q)q - \frac{1}{4}F_{\mu\nu}^a F_a^{\mu\nu}$$
$\bullet$ \textbf{Derivada Covariante:} $D_\mu = \partial_\mu - ig_s \left( \frac{\lambda ^a}{2}\right) A_\mu ^a$ \quad $g_s$ cte acoplamiento strong, \;  $a: 1-8$ \\
$\bullet$ \textbf{Tensor de Campo:} $F_{\mu\nu}^a = \partial_\mu A_\nu^a - \partial_\nu A_\mu^a + g_s f^{abc} A_\mu^b A_\nu^c$. \\
$\bullet$ \textbf{\hyperlink{su3}{Álgebra SU(3)}:}
$\setlength{\fboxsep}{1.5pt} \boxed{T^a=t^a = \tfrac{\lambda^a}{2}}$ \quad $[t^a, t^b] = i f^{abc} t^c$ \quad $f^{abc} = f_{abc}$ \quad $t^{a*} = t^a$.
\\
$\bullet$ \textbf{Invarianza gauge local:} $q(x) \rightarrow \exp{\left\{i\sum_a(\frac{\lambda^a}{2}) \theta_a(x)\right\}} q (x)$, ídem $D_\mu q(x)$.
\vspace{-5pt}
\begin{center}
\begin{tikzpicture}[scale=0.4, baseline=(current bounding box.center)]
    % Líneas de Fermiones (Quarks)
    \draw[fermion] (-1.5,-1) node[above left]{$q_j$} -- (0,0);
    \draw[fermion] (0,0) -- (1.5,-1) node[above right]{$q_i$};
    
    % Línea de Bosón (Gluón)
    \draw[gluon] (0,0) -- (0,1.5) node[above]{$A_\mu^a$};
    
    % Punto del Vértice
    \node[vertex] at (0,0) {};
    
    % Factor del Vértice (Etiqueta lateral)
    \node[right] at (1.95, 1.25) {$\displaystyle +ig_s \left( \frac{\lambda^a}{2} \right)_{ij} \gamma^\mu$}; 

    \node[right] at (1.95, 0.20) {$i,j$: índices de color};
\end{tikzpicture}
\end{center}
\vspace{-5pt}
\section{REGLAS DE FEYNMAN}
$\bullet$ \textbf{Propagadores:} Quark: $i\frac{(\cancel{p}+m)}{(p^2-m^2+i\epsilon)} \cdot \delta_{ij}$ \quad  \text{Gluón (Feynman):} $-i\eta_{\mu\nu}\frac{\delta^{ab}}{k^2}$. \\
$\bullet$ \textbf{Vértice Quark-Gluón:} $ig_s  (\frac{\lambda}{2}^a)_{ij}\gamma^\mu$. \\
$\bullet$ \textbf{Vértice 3-Gluones:} $-g_s f^{abc} [\eta_{\mu\nu}(k_1-k_2)_\rho + \eta_{\nu\rho}(k_2-k_3)_\mu + \eta_{\rho\mu}(k_3-k_1)_\nu]$. \\
$\bullet$ \textbf{Vértice 4-Gluones:} $-ig_s^2 [f^{abe}f^{cde}(\eta_{\mu\rho}\eta_{\nu\sigma}-\eta_{\mu\sigma}\eta_{\nu\rho}) + f^{ace}f^{bde}(\eta_{\mu\nu}\eta_{\rho\sigma}-\eta_{\mu\sigma}\eta_{\nu\rho}) + f^{ade}f^{bce}(\eta_{\mu\nu}\eta_{\rho\sigma}-\eta_{\mu\rho}\eta_{\nu\sigma})]$.

\section{CONSTANTES Y SUMAS DE COLOR (SU(3))}

$$C_F\!=\! \frac{N_c^2-1}{2N_c}\!=\!\frac{4}{3}\;\;\;\;C_A\!\!=\! N_c \!=\! 3\;\;\;\;T_F\!=\! \frac{1}{2}$$


$$\delta_{ab}C_A = \sum_{c,c'} f^{acc'}f^{bcc'}\quad \delta_{ij} C_F = \sum_{a,k}\frac{\lambda^a_{ik}}{2}\frac{\lambda^a_{kj}}{2}\quad\delta_{ab}T_F = \sum_{ij} \frac{\lambda^{a}_{ji}}{2} \frac{\lambda^{b}_{ij}}{2}$$



\section{HARD SCATTERING (LÍMITE ULTRARELATIVISTA $m \to 0$)}
Secciones eficaces promediadas ($\sum | \mathcal{M} |^2$) en QCD mediante comparación con QED: \\
$\bullet$ \textbf{$qq' \to qq'$ ($q \neq q'$):} Canal $t$. \quad 
$|\overline{\mathcal{M}}|^2 = \frac{2}{9} [2g_s^4 \frac{s^2+u^2}{t^2}]$. \\
$\bullet$ \textbf{$q\bar{q} \to q'\bar{q}'$ ($q \neq q'$):} Canal $s$.\quad
$|\overline{\mathcal{M}}|^2 = \frac{2}{9} [2g_s^4 \frac{t^2+u^2}{s^2}]$. \\
$\bullet$ \textbf{Bhabha ($q\bar{q} \to q\bar{q}$):} Canales $s$ y $t$. \quad $|\overline{\mathcal{M}}|^2 = 2g_s^4 [\frac{2}{9}\frac{u^2+t^2}{s^2} + \frac{2}{9}\frac{s^2+u^2}{t^2} - \frac{4}{27}\frac{u^2}{st}]$. \\
$\bullet$ \textbf{Compton QCD ($gq \to gq$):} Canales $s, u, t$ \quad
$|\overline{\mathcal{M}}|^2 = g_s^4 [\frac{s^2+u^2}{t^2} - \frac{4}{9}\frac{s^2+u^2}{su}]$. \\
$\bullet$ \textbf{$gg \to gg$:} $|\overline{\mathcal{M}}|^2 = \frac{9}{2} g_s^4 [3 - \frac{ut}{s^2} - \frac{st}{u^2} - \frac{su}{t^2}]$.

\section{CÁLCULO DE FACTORES DE COLOR ($F_c$)}
Procedimiento: extraer tensores de color de $\mathcal{M}$ y $\mathcal{M}^*$, promediar iniciales y sumar finales. \\
$\bullet$ \textbf{Promedio inicial:} $\frac{1}{N_{c1} N_{c2}}$ (Quarks: 3, Gluones: 8). \\
$\bullet$ \textbf{Ejemplo $qq' \to qq'$:} \quad $F_C =  \frac{1}{3 \cdot 3} \sum_{\substack{ijkm \\ a ba'b'}} (t^a_{ij}\delta_{ab}t_{km}^b)(t^{a'*}_{ij}\delta_{a'b'}t_{km}^{b'*})\!=\!\hspace{-5pt}\begin{array}{cc}
     \text{Hermiticidad}\\
      \text{+ props. $\delta_{ab}$}
\end{array} \hspace{-5pt}=$ \\
$= \frac{1}{9} \sum_{ijkm\,ab} (t_{ij}^a t_{km}^a)(t_{ji}^{a'}t_{mk}^{a'}) = \frac{1}{9} \sum_{a a'}\sum_{ij}(t_{ij}^a t_{ji}^{a'})\sum_{km}(t_{km}^a t_{mk}^{a'}) = \begin{array}{cc}
     \text{suma}\\
      \text{de color}
\end{array} = $\\
$=\frac{1}{9} \sum_{a a'}(T_F \delta_{aa'})(T_F \delta_{aa'})=\frac{1}{9} T_F^2 \sum_{a }\delta_{aa}^2 = \frac{1}{9}\cdot 8 \cdot T_F^2 = \frac{2}{9}$
\\
$\bullet$ \textbf{Patas con Quark y/o Fotones:} Incluir delta de color aunque el vértice no sea de QCD. \\
$\bullet$ \textbf{Procesos Mixtos:} $e^+e^- \to q\bar{q}$: $F_c = 3$ \quad  $\gamma q \to \gamma q$: $F_c = 1$ \quad $g q \to \gamma q$: $F_c = \frac{1}{6}$.

\section{QCD EN COLISIONADORES $e^+e^-$ Y COLISIONADORES HADRÓNICOS}
Los partones producidos en el estado final se confinan formando \textbf{jets}. La configuración cinemática típica en colisión frontal cumple $p_{j_1} + p_{j_2} + \dots = 0$.

\textbf{Producción de 2 jets}: $\sigma$($e^-e^+ \to$ 2 jets) $\approx \sum_q \sigma(e^+e^- \to q \bar{q})$. Mediante comparación con $e^+e^- \to \mu^- \mu ^+$, se obtiene para $e^+e^- \to q \bar{q}$:
$$|\overline{F}|^2 = e^4 Q_q^2 N_c \frac{t^2+u^2}{s^2}$$
 $$\sigma(e^-\!e^+\! \to \!q\bar{q}) \!=\! \frac{4\pi\alpha^2}{3s} (Q_q^2 N_c) \!=\! \sigma(e^-\!e^+\! \to\! \mu^-\!\mu^+\!)(Q_q^2 N_c) \quad \frac{\sigma(e^-\!e^+\! \to\! \text{2 jets})}{\sigma(e^-\!e^+\! \to\! \mu^-\!\mu^+)} \!=\! N_c\! \sum_q Q_q^2$$
\\ $\bullet$ $\sum_q$ depende de $\sqrt{s} = 2E$:\;  $\sqrt{s} = [5, 20, 10^3]$ GeV $\to q: [u\,d\,s\,c,\; u\,d\,s\,c\,b, \; u\,d\,s\,c\,b\,t]$


\textbf{Producción de 3 jets:} ($e^-e^+ \to$ 3 jets) $\approx \sum_q \sigma(e^-e^+ \to q\bar{q}g)$. El factor de color es constante para las contribuciones.
$F_c = \sum_{\alpha, i, j} t_{ij}^\alpha t_{ji}^\alpha = C_F \cdot N_c = \frac{4}{3} \cdot 3 = 4$

Definiendo las fracciones de energía del centro de masas para los partones finales $x_i = \frac{2E_i}{\sqrt{s}}$ sujeto a la restricción cinemática $x_1+x_2+x_3=2$, la ratio diferencial es:
$$\frac{\sigma(e^-e^+ \to 3\text{ jets})}{\sigma(e^-e^+ \to 2\text{ jets})} = \frac{2\alpha_s}{3\pi} \iint \frac{x_1^2+x_2^2}{(1-x_1)(1-x_2)} dx_1 dx_2\qquad \alpha_s = \frac{g_s^2}{4\pi}$$

\textbf{Colisón de hadrones ($pp$ o $p\bar{p}$):}  los partones no son elementales. Intervienen \underline{\text{PDFs}} en el estado inicial y \underline{\text{funciones de hadronización}} en el final.

$\bullet$ \text{Teorema de Factorización:}
La sección eficaz hadrónica total es la convolución de las PDFs con la sección eficaz partónica (\textit{hard scattering process}):
$$\sigma(pp \to X) = \sum_{i,j} \int_0^1 dx_1 \int_0^1 dx_2 \, f_i(x_1, Q^2) f_j(x_2, Q^2) \hat{\sigma}_{ij}(\hat{s}) + (x_1 \leftrightarrow x_2)$$
\\ $\bullet$ $f_{i/j}(x, Q^2)$: prob hallar partón $i/j$ con fracción de momento $x$ evaluado a energía $Q^2$.
\\ $\bullet$ $\hat{\sigma}_{ij}$: Sección eficaz a nivel de partones (ej. $q\bar{q} \to q'\bar{q}'$). Dominado por QCD: $\hat{\sigma} \sim \mathcal{O}(\alpha_s^2)$.
\\ $\bullet$ $\hat{s}= x_1 x_2 s$: Invariante de Mandelstam partónico. Se asume típicamente $Q^2 = \hat{s}$.
\\ $\bullet$ Suma contribuciones a nivel de sección eficaz (no de amplitud, $\nexists$ interferencia macro).\\
$\bullet$ \textbf{Aproximación Valencia:} Hadrón $H$ compuesto por $n_v$ quarks de valencia, PDF para quark $i$: 
$q_{v,i}(x, Q^2) = n_i \delta\left(x - \frac{1}{n_v}\right)$, con $n_i$ quarks de valencia de sabor $i$ en $H$.

